\chapter{Introduction}

\section{Background and Motivation}

Cloud computing adoption has accelerated significantly with major vendors such as Amazon Web Services (AWS) and Microsoft Azure continually expanding proprietary service offerings. These cloud platforms provide a wide range of infrastructure, platform, and software services catering to diverse business needs. However, the growing variety and complexity of cloud offerings introduce significant challenges for solution architects who must balance competing factors such as cost, scalability, security, and performance when designing cloud architectures. %\cite{agiliway2025aws-vs-azure}.

The present study proposes an innovative agentic AI framework to assist architects in navigating the resulting architectural complexity. The framework employs an Evaluator-Optimizer workflow \cite{anthropic2024agents} that iteratively generates, validates, and refines cloud architectural recommendations. Validation is grounded in authoritative cloud provider documentation through Retrieval Augmented Generation (RAG) \cite{lewis2020rag, asai2023selfrag}, ensuring factual and reliable outputs. The system aims to deliver efficient, secure, and cost-optimized solutions capable of cross-provider comparative analysis. %\cite{futransolutions2025multi-cloud}.

\section{Problem Statement}

The transition from traditional on-premises infrastructure to cloud platforms has resulted in an exponential increase in cloud services. Providers like AWS and Azure often offer overlapping services solving similar problems, complicating architectural decisions. %\cite{simplilearn2025.aws-vs-azure-challenges}
Architects typically rely on domain expertise and experience to evaluate tradeoffs involving costs, scalability, security, and performance. Current AI-based tools and assistants, including first-party vendor AI products like Amazon Q \cite{aws2023q} and Microsoft Copilot \cite{microsoft2024copilot}, suffer from several limitations:

\begin{itemize}
    \item Existing chatbots often provide partial guidance focusing on isolated components rather than comprehensive end-to-end architectures.
    \item Generated outputs frequently lack grounded factual validation against authoritative cloud documentation.
    \item Current solutions generally do not support comparative evaluation across multiple cloud providers, often leading to vendor-biased recommendations.
\end{itemize}

The outlined challenges highlight the need for an intelligent AI-driven framework capable of generating, validating, and optimizing architectural recommendations with multi-cloud comparative features to aid in unbiased, informed decision-making. %\cite{agiliway2025aws-vs-azure, futransolutions2025multi-cloud}.

\section{Objectives}

The primary objectives of this research project are:

\begin{itemize}
    \item To design and implement an Evaluator-Optimizer feedback workflow using Retrieval Augmented Generation (RAG) for factually grounded recommendation validation.
    \item To incorporate cost and feature tradeoff analysis to enhance efficiency, security, and cost optimization in architectural decisions.
    \item To develop autonomous AI agent teams representing AWS and Azure cloud providers capable of generating domain-specific cloud architecture proposals.
    \item To enable multi-turn debate mechanisms facilitating side-by-side comparative evaluation of architectures proposed by different agent teams.
    \item To design and deploy a functional prototype system accessible through simple user interfaces.
    \item To validate the effectiveness of the framework via measurable performance metrics, benchmark case studies, and domain expert evaluations.
\end{itemize}

\section{Scope and Limitations}

\subsection{Scope}

The research focuses on developing and evaluating a functional prototype of the agentic AI framework, limited to the following boundaries:

\begin{itemize}
    \item \textbf{Cloud Providers:} Initially limited to Amazon Web Services (AWS) and Microsoft Azure, the two dominant public cloud vendors.
    \item \textbf{Knowledge Source:} The Retrieval Augmented Generation knowledge base is constructed exclusively from publicly available, official provider documentation such as whitepapers, service guides, and recommended best practices.
    \item \textbf{Output Format:} The system outputs comprehensive textual cloud-architecture recommendations with detailed justifications, as well as generates deployable code (Infrastructure-as-Code templates such as Terraform or CloudFormation).
\end{itemize}

\subsection{Limitations}

The implemented prototype currently has the following operational constraints:

\begin{itemize}
    \item \textbf{Absence of Real-Time Data:} The framework does not incorporate real-time information such as pricing updates, service availability in specific regions, or customer-specific account configurations.
    \item \textbf{Knowledge Base Boundaries:} Although extensive, the knowledge base may lack coverage of some niche or recently introduced cloud services due to document availability and constraints associated with the open source vector database.
    \item \textbf{Non-Substitutive Advisory Role:} The system is intended as a decision support tool complementing human expertise rather than replacing human operators. All architectural recommendations warrant expert review prior to deployment.
\end{itemize}

\section{Feasibility Study}

\subsection{Technical Feasibility}

Mature underlying technologies such as Large Language Models (LLMs), vector similarity searches, and agentic AI frameworks (e.g., LangChain \cite{langchain2024docs}, LangGraph \cite{langgraph2024docs}, AutoGen \cite{wu2023autogen}) enable the realization of the proposed framework. Cloud platforms (AWS and Azure) provide rich documentation and resources essential for building the knowledge base.

\subsection{Data Availability}

Cloud vendor documentation required for the RAG knowledge base is openly accessible. Such resources include service whitepapers, architecture best practices, and official guides, forming a comprehensive foundation for reliable information retrieval and validation. %\cite{aws2025well-architected, azure2025-architecture}.

\subsection{Operational Feasibility}

The prototype targets deployment on affordable local computation platforms using GPU-accelerated frameworks (Apple Metal Performance Shaders \cite{apple2024mps}) and local embeddings (Ollama \cite{ollama2024web}), with the potential to integrate cloud based services. The user interfaces will be minimal yet practical comprising command-line and web portals, to streamline accessibility, including for evaluation purposes.

\subsection{Expertise and Challenges}

The project benefits from the author's background in data engineering, AI/ML, cloud computing, and Python development. The primary anticipated challenge is ensuring high accuracy and factual reliability of AI-generated recommendations, requiring robust knowledge ingestion and validation techniques.

\section{Thesis Organization}

The dissertation is organized into the following chapters:

\begin{itemize}
    \item \textbf{Chapter 2} provides an extensive literature review, outlining foundational concepts in agentic systems, RAG methodologies, and related areas, highlighting existing challenges.
    \item \textbf{Chapter 3} details the system design and methodology, explaining the agentic AI architecture, workflow orchestration, and implementation specifics.
    \item \textbf{Chapter 4} details the core implementation components of the agentic AI framework, including web scraping, vector database construction, and the multi-agent workflow.
    \item \textbf{Chapter 5} presents the progress and findings from the first phase of the project, including project status, preliminary testing, and quantitative evaluation results.
    \item \textbf{Chapter 6} concludes the thesis with summaries, limitations, and prospects for future work, including roadmap extensions and multi-cloud comparison capabilities.
\end{itemize}